\section{Method}
\label{sec:method}
\subsection{Predict: a local map from two images}
\label{sec:predict}
A current image $I_f$ and keyframe image $I_k$ enter one shared, frozen
network. Its pointmaps and matching features feed the MASt3R-SLAM tracker.
Its Gaussian head predicts a primitive at each valid pixel, in the first
camera's coordinates. A primitive has centre $\mu_n$, covariance $\Sigma_n$,
opacity $\alpha_n$, and spherical-harmonic (SH) colour. For predicted point
$x_n$, offset $\Delta_n$, positive scale $s_n$, and unit quaternion $q_n$,
\begin{equation}
  \mu_n=x_n+\Delta_n,\qquad
  \Sigma_n=R(q_n)\operatorname{diag}(s_n^2)R(q_n)^\top ,
  \label{eq:gaussian}
\end{equation}
The DC colour coefficient is a residual added to the SH encoding of input
RGB; conversion clamps the resulting colour.
The tracker maintains keyframe poses
$T_k=(a_kR_k,t_k)\in\mathrm{Sim}(3)$, where $a_k>0$ allows local scale
correction. Matching, pointmap fusion, and pose-graph optimisation follow
MASt3R-SLAM. The mapping path reads these poses.

\subsection{Anchor: move the map and its observations together}
\label{sec:anchor}
Each keyframe stores its Gaussians in local coordinates. Rendering places
them in the world with the current anchor pose:
\begin{equation}
  \mu_n^W=a_kR_k\mu_n^k+t_k,\qquad
  \Sigma_n^W=a_k^2R_k\Sigma_n^kR_k^\top .
  \label{eq:anchor}
\end{equation}
A loop closure can then move a local map by updating one pose.
The observation cameras must follow the same correction. For a tracked
frame $f$ associated with keyframe $k$, we store its image and relative
transform $T_{kf}$. Its current world pose is
\begin{equation}
  T_f^W=T_k^W T_{kf}.
  \label{eq:supervision}
\end{equation}
For this camera and a local map sharing anchor $k$, the anchor cancels:
\begin{equation}
  (T_f^W)^{-1}T_k^W=(T_k^WT_{kf})^{-1}T_k^W=T_{kf}^{-1}.
  \label{eq:consistency}
\end{equation}
Thus their relative coordinates survive any invertible update of $T_k^W$.
A fixed world-space supervision camera would break this relation.
Equation~(\ref{eq:consistency}) concerns one shared anchor; overlapping maps
with different anchors can still disagree.

\subsection{Refine: fit the assembled map to tracked images}
\label{sec:refine}
The refiner renders the current map and compares it with sampled RGB
observations. It optimises log-scales, rotations, opacity logits, and DC
colour with
\begin{equation}
  \mathcal L_{\mathrm{ref}}
  =0.8\,\|\hat I-I\|_1+
    0.2\,\bigl(1-\operatorname{SSIM}(\hat I,I)\bigr),
  \label{eq:ref}
\end{equation}
where the norm is mean absolute pixel error and SSIM follows
\cite{wang2004ssim}. Gaussian centres are fixed in the external comparison.
Adam~\cite{kingma2015adam} uses learning rates $2.5{\times}10^{-3}$ for DC
colour, $5{\times}10^{-2}$ for opacity, $5{\times}10^{-3}$ for scale, and
$10^{-3}$ for rotation.

The supervision pool has a reservoir of 200 historical frames and a ring of
64 recent frames. Sampling selects the recent ring with probability 0.3.
Tracked non-keyframes can enter the pool. Each step reconstructs their
cameras with (\ref{eq:supervision}), then updates the Gaussian parameters.
The tracking model and camera poses receive no photometric gradients.

Refinement runs in a separate process, optionally on a second GPU. New
Gaussian records transfer once; current anchor poses are read at each step.
The refiner publishes complete, versioned CPU snapshots for display.
On one GPU, a duty cycle limits competition with tracking. An optional
post-sequence \emph{polish} phase continues with the final pose graph.

\subsection{Improve the starting point}
\label{sec:initialisation}
\paragraph{Opacity at insertion}
Different keyframes can predict the same surface. Their contributions
interact when the renderer composites them. For projected opacity $b_n(u)$,
\begin{equation}
  \hat I(u)=\sum_n c_n(u)b_n(u)\prod_{j<n}(1-b_j(u)).
  \label{eq:compositing}
\end{equation}
Reducing opacity can reduce one layer's contribution and reveal layers
behind it. We test this adjustment at map insertion.
Let $\hat c_n\in[0,1]$ be the rank-normalised pointmap confidence within
one keyframe. We initialise opacity as
\begin{equation}
  \alpha'_n=\alpha_n
    \operatorname{clip}\!\left(1-2\lambda(1-\hat c_n),\,0.1,\,1\right).
  \label{eq:thin}
\end{equation}
Ranks avoid a calibration assumption across frames. The experiments compare
$\lambda=0.45$ with $\lambda=0$ and with uniform attenuation.
Every primitive remains available for later optimisation because the
multiplier is at least 0.1.

\paragraph{Head adaptation}
We also test whether family-specific training improves the predicted map.
Only the Gaussian DPT head is trainable; parameter hashes verify that the
shared encoder and decoder remain unchanged. The upstream objective is
\begin{equation}
  \begin{aligned}
  \mathcal L_{\mathrm{head}}={}&
  \mathcal L_{\mathrm{MSE}}(M\odot\hat I,M\odot I)\\
  &+0.25\,\mathcal L_{\mathrm{LPIPS}}(M\odot\hat I,M\odot I),
  \end{aligned}
  \label{eq:head}
\end{equation}
where $M$ selects target pixels visible from the context pair.
Depth or pseudo-depth is used to construct samples and visibility masks.
This loss fits rendered appearance without an explicit geometric error term.
Head adaptation changes prediction, opacity attenuation changes insertion,
and refinement fits the assembled sequence. The experiments isolate these
three stages.
